\chapter{Introduction}
\label{chapter1:introduction}
\section{Introduction}

Modern high-energy physics seeks to understand the elementary constituents of matter and the fundamental interactions that govern their behaviour. 
At present, the most successful theoretical description of particle physics is the Standard Model~(SM), a quantum field theory based on the gauge symmetry
\begin{equation}
    SU(3)_{C} \times SU(2)_{L} \times U(1)_{Y}.
\end{equation}
Within this framework, matter is described by three generations of quarks and leptons, while the electromagnetic, weak, and strong interactions are mediated by gauge bosons. 
The electroweak symmetry is spontaneously broken through the Higgs mechanism, giving masses to the weak gauge bosons and to the fermions through their interactions with the Higgs field. 
The discovery of the Higgs boson by the ATLAS and CMS experiments in 2012 completed the particle content predicted by the SM and marked a major milestone in modern particle physics.

The SM has been tested with remarkable precision over a wide range of energy scales and has successfully described a vast number of experimental observations. 
Nevertheless, it is widely regarded as an incomplete theory. 
It does not provide a particle candidate for dark matter, does not explain the observed matter--antimatter asymmetry of the universe, and does not account for the origin of non-zero neutrino masses. 
In addition, the pattern of fermion masses and flavour mixing remains unexplained, and the large separation between the electroweak and Planck scales suggests the possible existence of new physics beyond the SM. 
These open questions motivate direct searches for new particles and interactions at the highest experimentally accessible energies.

The Large Hadron Collider~(LHC) at CERN provides a unique environment for probing such new phenomena. 
By colliding protons at tera-electronvolt centre-of-mass energies, the LHC allows the production of particles that may be associated with physics beyond the SM. 
During Run~2 of the LHC from year 2015 to 2018, proton--proton~($pp$) collisions were delivered at a centre-of-mass energy of $\sqrt{s} = 13~\TeV$.

ATLAS is a detector located at one of the interaction points of the LHC. 
It is designed to observe and study a wide range of physics processes, mainly from the $pp$ collisions, from precision measurements of SM interactions to searches for new particles. 
The detector consists of several sub-detectors surrounding the collision interaction point, including an inner tracking detector, electromagnetic and hadronic calorimeters, and a muon spectrometer. 
Together with a powerful magnet system, these components allow the reconstruction and identification of electrons, photons, muons, hadronically decaying $\tau$-leptons, jets, and missing transverse momentum, denoted by $\met$.

Among the many possible extensions of the SM, leptoquarks~(LQs) provide a well-motivated class of new particles connecting the quark and lepton sectors. 
They are hypothetical bosons that couple directly to a quark and a lepton, and they appear naturally in theories such as grand unified theories and models of quark--lepton unification. 
Leptoquarks are also of particular phenomenological interest because they can modify semileptonic processes and may therefore be connected to the observed tensions in flavour observables.

A search for the leptoquarks is presented in the thesis using 140~\ifb of $pp$ collision data recorded by the ATLAS detector at a centre-of-mass energy of $\sqrt{s}=13~\TeV$. 
This thesis focuses on the vector singlet leptoquark $U_{1}$ as the benchmark model. 
This model can contribute to processes involving $b$ quarks, $\tau$ leptons, and neutrinos.
Therefore, the analysis targets events with a hadronically decaying $\tau$-lepton, large missing transverse momentum, and jets. 
The observed data are compared with SM background predictions in signal regions designed to enhance sensitivity to possible leptoquark contributions. 
In the absence of a significant excess over the SM expectation, constraints are placed on the allowed parameter space of the $U_{1}$ model, expressed in terms of the leptoquark mass and its couplings.

This dissertation is organised as follows. 
Chapter 2 provides an overview of the SM's theoretical framework and the physics beyond it.
Chapter 3 describes the leptoquark models and the searches performed at LHC. 
Chapter 4 introduces the LHC and the ATLAS detector.
Chapter 5 provides the dataset and simulated samples.
Chapter 6 details the reconstruction and the identification of the physics objects that are used for event selection.
Chapter 7 discusses the optimisation for the event selection used for signal signature extraction.
Chapter 8 presents the background estimation strategy.
Chapter 9 discusses the systematic uncertainties for the background and signals.
Chapter 10 presents the statistical interpretation of the results.
Chapter 11 summarises this analysis.
